\section{Related Work}
\label{sec:related}

\paragraph{Manually-designed kernels.} Many existing frameworks, such as TensorFlow XLA~\cite{tensorflow_xla, Tensorflow}, PyTorch~\cite{pytorch}, and TensorRT~\cite{tensorrt}, rely on GPU experts to manually design kernels for ML operators.
%
Recently, significant engineering effort has been dedicated to hand-optimizing GPU kernels for commonly used DNNs, particularly foundation models~\cite{bommasani2022opportunities}.
For example, to accelerate attention computation~\cite{transformers}, several specialized kernels have been developed based on FlashAttention~\cite{dao2023flash, tri2023flashdecoding, hong2024flashdecoding, fastertransformer}. 
Due to the increasing complexity of modern GPUs---such as tensor cores in A100s~\cite{Markidis2018tensorcores} and thread block clusters in H100s~\cite{nvidia-h100}---manually designed kernels may miss subtle optimizations that are hard to discover manually.


\paragraph{Superoptimization-based approaches.} Superoptimization was originally introduced to find optimal instruction sequences~\cite{massalin, stoke, peephole}. 
%
Recent work has applied superoptimization techniques to tensor programs~\cite{TASO, wang2021pet, zheng2023einnet, tensat, unger2022unity, FlexFlow, hu2024korch, jeon2025graphpipe}.
%
However, all these attempts only consider algebraic transformations at the kernel level and cannot discover more sophisticated optimizations that require jointly considering algebraic and schedule transformations at all of the kernel, block, and thread levels.
Our evaluation shows that \sys largely outperforms existing DNN superoptimizers, demonstrating the importance of multi-level joint optimization.

\paragraph{Schedule-based approaches.} Recent work has introduced ML compilers that automatically optimize the execution schedule of kernel GPUs. Systems such as TVM~\cite{tvm, tvm_auto_tuner}, Ansor~\cite{ansor}, and Triton~\cite{tillet2019triton}, along with others~\cite{zheng2020flextensor, hagedorn2023graphene, feng2022tensorir}, build on the idea of algorithm-schedule separation introduced in Halide. They search for optimized schedules to execute a user-specified algorithm on GPUs.
However, schedule-based approaches require users to explicitly specify the algorithm for each kernel, and their performance is limited to the quality of these provided algorithms.

\paragraph{Multi-level graph representations.} Welder~\cite{shi2023welder} and ASPEN~\cite{park2023aspen} introduce multi-level tile graphs that share similarities with \sys's \graphs, as both representations follow the GPU hierarchy. However, prior work focuses on scheduling transformations, while \sys extends beyond scheduling by also considering algebraic transformations and the discovery of new custom kernels. Most optimizations presented in this paper fall outside the scope of these prior approaches.

% A key difference between \sys and existing automated approaches is that \sys does not require users to manually specify the computation of a kernel. Instead, \sys performs superoptimization at the kernel, block, and thread levels and jointly considers algebraic and schedule transformations at all these levels to discover currently missing optimizations.